% !TeX root = ../papers/nc-p3.tex
\section{正則性平面の小標数の全有理点}
\label{sec:rank-fifteen-quotient-residues}
\begin{theorem}\label{thm:rank-fifteen-quotient-residues}
剰余標数が $3,5,7$ のいずれかの付値を考える。標数零の通常な階数15窓の
正則性座標 $x,y$ が整で、$y$ は単元、両者の剰余が素体に属するとする。
この窓は標準Koszul左右AS正則であり、四つの中心四次式の多項式環上で
各始点の経路の直和が階数256の自由加群となる。
倍率や四次根自体の剰余には素体条件を課さない。
\end{theorem}
\begin{proof}
$v=V/(U-1)^2$ とすると正則性平面からの一つの持上げは
\[
 v(x-4v)^2=y,\qquad U=x/v-3,\qquad V=y/v^2
\]
により計算できる。$2,y$ は単元なので、三次式の最高係数と定数項は単元である。
従って $v,x-4v,U-1,V$ は全て単元である。次に $K^2=U,H^2=V$、
$r+1/r=2(K+H),w+1/w=2(K-H)$、$s^2=r,t^2=w$ と
$n^4+(1-U)n^2+V=0$ の根を取る。$s,t,n$ とその逆は整である。
これらは順にモニックな方程式なので、代数閉体の付値を延長すれば指定した剰余根に
落ちる根が各段で少なくとも一つある。三次式や四次式の還元が重根を持つ場合も同じである。

\path{computations/verify_rank_fifteen_quotient_residues.py} は各剰余標数の
全ての $(\bar x,\bar y)\in\mathbb F_p\times\mathbb F_p^\times$ を調べ、
上の方程式を有限拡大体で順に解く。検証した点は順に6,20,42個である。必要になった拡大次数は2,4,6,12で、
各次数の証明書の個数は標数3で $(1,3,1,1)$、標数5で $(2,10,6,2)$、
標数7で $(5,21,12,4)$ となる。
選んだ一つの $(s,t,n)$ から規格化テンソルを作り直し、その隣接像で
四始点の関係を生成する。隣接階数6、長さ4,5の経路次元 $35,56$、
中心の文字次元 $(3,2,0,0)$ を確認し、四つの中心式による商の長さ0から13の
次元を全始点で検査する。例えば $p=7,(\bar x,\bar y)=(0,2)$ では
文字零の中心の三自由座標上で $(1,0,1),(0,1,1)$ を使い所定の次元列を得る。

同じ標数零の $(x,y)$ の元の通常窓と、選んだ形式的な持上げのテンソルは
系\ref{cor:rank-fifteen-regularity-quotient}の証明により、基底変更・巡回移動・反転で対応する。
従って持上げ先の隣接像には経路と中心の所定の次元の下界がある。
剰余の階数と四中心商の証明書が定理\ref{thm:rank-fifteen-regular-orbits}と
同じ方法で持ち上がり、元の窓の全次数の主張が従う。
\end{proof}
